\chapter{CONCLUSIONES Y RECOMENDACIONES}

\section{CONCLUSIONES}

Se implementó un sistema web que apoya la organización y administración del Desafío Bebras en Bolivia: desde la información pública y la inscripción de maestros y estudiantes hasta la rendición, los resultados y los certificados. Respecto de cada objetivo específico:

\begin{itemize}
    \item \textbf{Gestión de contenido.} El CMS permite editar las páginas del sitio informativo armándolas con bloques, administrar su menú y su pie, crear páginas nuevas, publicar noticias e imágenes, guardar y recuperar respaldos, y publicar los cambios, al instante o en una fecha programada. Sus 254 pruebas aprueban.
    \item \textbf{Gestión de desafíos.} El banco de tareas representa las 43 tareas del cuadernillo Bebras 2024 con seis tipos de respuesta, cada una con su dificultad por categoría de edad y su probador. Un desafío reúne varias categorías con sus preguntas, sus fechas y su puntaje, y avanza solo por sus etapas. La práctica pública permite ejercitarse sin cuenta.
    \item \textbf{Registro y administración de participantes.} Los maestros se registran con su unidad educativa del catálogo nacional y el organizador los aprueba con una carta del director. Cada maestro inscribe a sus estudiantes en grupos por categoría, uno a uno, con planilla o compartiendo un código, y cada estudiante rinde con un código personal, sin necesidad de cuenta.
    \item \textbf{Evaluación.} El estudiante rinde en línea con sus respuestas guardadas y corregidas al instante; los puntajes, el ranking por categoría y los resultados se publican solos en la fecha fijada. La plataforma soportó 1600 estudiantes a la vez sin errores y con todos los puntajes exactos. Para las aulas sin conectividad, el maestro imprime la prueba y carga después las hojas de respuestas, que se corrigen como las demás.
\end{itemize}

Las pruebas automatizadas acompañaron cada iteración, como pide XP, y encontraron defectos que habrían impedido un desafío nacional antes de que llegaran a los usuarios. El traslado de la plataforma a Cloudflare hizo que el sistema escale con la demanda y se despliegue solo desde el repositorio.

\section{RECOMENDACIONES}

\begin{itemize}
    \item Incorporar la prueba en papel a la batería de pruebas automatizadas.
    \item Realizar pruebas de usabilidad con estudiantes y maestros reales antes de la primera edición.
    \item Repetir la prueba de estrés en el entorno de pruebas de Cloudflare, para medir los tiempos reales.
    \item Mejorar el rendimiento de la portada del sitio informativo, reduciendo el peso de sus imágenes.
\end{itemize}
